\documentclass[10pt,a4paper]{amsart}
\usepackage[margin=19mm]{geometry}
\usepackage{amsmath,amssymb,mathtools}
\usepackage[scr=rsfs,cal=euler]{mathalfa}
\usepackage{xfrac}
\usepackage[unicode,hidelinks,bookmarks=false]{hyperref}
\newcommand{\ie}{i.e.\ }
\newcommand{\cf}{cf.\ }
\newcommand{\resp}{respectively\ }
\newcommand{\Subsection}{\subsection}
\providecommand{\nobreakdash}{\nobreak\hbox{-}}
\setlength{\emergencystretch}{2em}
\multlinegap0pt
\usepackage{amssymb}
\usepackage{xcolor}
\newtheorem{lemma}{Lemma}
\newtheorem{proposition}{Proposition}
\def\a{\alpha}
\def\b{\beta}
\def\lr{\leftrightarrow}
\title{\large\textsc{Hatcher--Thurston complex for surfaces with non-planar ends}}
\author{Manvendra Somvanshi}
\date{Source: arXiv:2602.22393v1 (CC BY 4.0)}
\begin{document}
\maketitle

\noindent\emph{Source and licence.} \large\textsc{Hatcher--Thurston complex for surfaces with non-planar ends}. Version arXiv:2602.22393v1 (CC BY 4.0), distributed by arXiv under the Creative Commons Attribution 4.0 International licence, http://creativecommons.org/licenses/by/4.0/, as recorded in the arXiv licence metadata for this exact version and shown on the abstract page https://arxiv.org/abs/2602.22393v1. Full licence notice: \texttt{licenses/arxiv-2602.22393-CC-BY-4.0.txt}.\par\smallskip

\noindent\textit{Editorial scope.} Counted unit: the article's proposition prop:base-case (source0.tex lines 266--268). The article's complete original proof of it is delivered with it (source0.tex lines 269--274, one \texttt{proof} environment of 6 lines). No mathematics is written, repaired or re-derived: the statement and the proof are the article's own lines, in the article's own notation, and no retained line is substituted or re-flowed.  Retained uncounted as the source-local context the proof needs: lines 209--211.  Nothing was omitted from any retained span, so the package carries no omission record.  No retained line contains a citation, so the wrapper carries no bibliography and no bibliography entry.  No retained line contains a figure environment, an image-inclusion command or drawing code, so no image is delivered and the package declares no asset dependency.  Editorial formatting only, in addition: an AMS \texttt{amsart} preamble that restates the article's own declarations and macros used by the retained lines, namely the environments \texttt{lemma}, \texttt{proposition} on separate counters, together with the standard packages those lines need.  The retained environments print in this document as Lemma 1, Proposition 1 in order of appearance, whereas the article prints the same items with the numbering of its own full document.  The complete native source of the article as distributed in the arXiv e-print for this exact version is bundled unchanged as \texttt{sources/195329/source0.tex}.
\begin{lemma}\label{lem:radius-1}
  Every closed path $p=a_0\lr \cdots \lr a_{n-1} \lr a_0$ in $\Gamma_1(S)$ of radius $1$ with respect to $a_0$ is null-homotopic. 
\end{lemma}

\begin{proposition}\label{prop:base-case}
  Every closed path $p=a_0\lr\cdots\lr a_{n-1}\lr a_0$ in $\Gamma_1(S)$ is null-homotopic.
\end{proposition}

\begin{proof}
  The idea of the proof is to construct curves $\b_0,\ \b_1,\ \b_2$ in $S$ so that $i(b_2,a_j) = 0$ for all $j$ and $i(b_2,b_0) = i(b_2,b_1)= i(b_0,a_0) = i(b_1,a_1) = 1$. This means that the path $q= b_2 \lr b_1 \lr a_1 \lr \cdots \lr a_{n-1}\lr a_0 \lr b_0 \lr b_2$ has radius 1 with respect to $b_2$. Since the path $b_2\lr b_1\lr a_1 \lr a_0\lr b_0\lr b_2$ is also radius $1$ with respect to $b_2$ it follows from Lemma \ref{lem:radius-1} that $q$ is homotopic to $p$ and that $q$ is null homotopic. It follows that $p$ is null-homotopic.\\


  Now we describe the construction of the curves $\b_i$. Let $R$ be a subsurface of $S$ with compact boundary and at least one non-planar end so that all the curves $\a_0,\cdots, \a_{n-1}$ are contained in $R$. Let $\delta$ be a separating curve in $R$ disjoint from all $\a_i$ so that the two surfaces, $R_1$ and $R_2$, obtained by cutting $R$ along $\delta$ satisfy the following: $R_1$ is compact and contains all $\a_i$ and $R_2$ is of infinite-type with a single compact boundary component corresponding to the curve $\delta$ in $R$ and has non-planar ends. Let $\tau_0$ and $\tau_1$ be homotopically distinct arcs in $R_2$ with their end points lying on the boundary of $R_2$ so that $\tau_0\cap \tau_1 = \emptyset$. Since $\delta$ and $\a_0$ are disjoint and $i(a_0,a_1)=1$, it can be shown that there exist disjoint essential arcs, $\sigma_0,\ \sigma_1$, with end points on the boundary of $R_1$ corresponding to $\delta$ and $|\sigma_0\cap \a_0| = 1$ and $|\sigma_1\cap \a_1|=1$. Gluing $R_1$ and $R_2$ together along the boundary corresponding to $\delta$ in such a way that $\beta_0 = \sigma_0\cup \tau_0$ and $\beta_1 = \sigma_1\cup \tau_1$ are closed curves. Then by construction we have $i(b_0,a_0) = i(b_1,a_1) = 1$ and that $b_i$ are non-separating. Let $\b_2$ be a non-separating curve in $R_2$ which intersects $\tau_0$ and $\tau_1$ exactly once. Thus we have our desired curves $b_0,\ b_1,$ and $b_2$.
\end{proof}

\end{document}
